\begin{table}[H]
\centering
\caption{外部相关系统与本文实验的关系}
\label{tab:external_references}
\small
\begin{tabularx}{\linewidth}{>{\raggedright\arraybackslash}p{0.23\linewidth} >{\raggedright\arraybackslash}p{0.20\linewidth} >{\raggedright\arraybackslash}X >{\raggedright\arraybackslash}X}
\toprule
系统/工作 & 任务重点 & 报告口径 & 与本文关系 \\
\midrule
Onsets and Frames \cite{hawthorne2018onsets} & 钢琴 AMT & onset/frame 联合转录 & 为 onset-gated 解码提供方法基础。 \\
ByteDance high-res. \cite{kong2021highres} & 钢琴 AMT 与踏板 & MAESTRO onset/offset/pedal & 作为高分辨率 onset/offset 与踏板建模的强参考。 \\
MT3 \cite{gardner2022mt3} & 多任务多音轨 AMT & 序列到序列转录 & 作为早期 seq2seq 探索与多任务建模参考。 \\
hFT-Transformer \cite{toyama2023hft} & 钢琴 AMT & 频率--时间层级建模 & 说明钢琴 AMT 仍可通过声学结构建模提升。 \\
Basic Pitch \cite{basicpitch2022} & 通用音频转 MIDI & 轻量开放工具 & 作为通用转录工具的外部参照，不直接比较数值。 \\
PM2S \cite{liu2022pm2s} & performance MIDI 到乐谱 & beat tracking 与 score conversion & 支撑谱面层独立建模的必要性。 \\
Score Transformer \cite{suzuki2021scoretransformer} & note-level 到 score & 谱面 token 生成 & 说明 MIDI 事件之外仍需预测记谱属性。 \\
MIDI-to-score Transformer \cite{beyer2024midi2score} & performance MIDI 到 score & 端到端序列转换 & 作为学习式谱面层的参考方向。 \\
\bottomrule
\end{tabularx}
\end{table}
